% Copyright (c) 2026 Eugene Joseph M. Ragasa.
% SPDX-License-Identifier: Apache-2.0
% Author: Eugene Joseph M. Ragasa
% Formal proof completed: 2026-10-29
\documentclass[11pt]{article}
\usepackage[T1]{fontenc}
\usepackage{amsmath,amssymb,booktabs,tabularx}
\usepackage[margin=1in]{geometry}
\usepackage[hidelinks]{hyperref}
\usepackage{listings}
\lstset{basicstyle=\ttfamily\small,breaklines=true,columns=fullflexible,
  literate={∃}{{$\exists$}}1 {∈}{{$\in$}}1 {↔}{{$\leftrightarrow$}}1
           {∩}{{$\cap$}}1 {¬}{{$\neg$}}1}
\newcolumntype{Y}{>{\raggedright\arraybackslash}X}
\title{Proof Dossier: Compact Positive Distance and Disjointness}
\author{Eugene Joseph M. Ragasa}
\date{Formal proof completed: 2026-10-29}
\begin{document}\maketitle
\section{Identity and status}
\begin{tabularx}{\textwidth}{@{}lY@{}}\toprule
Statement ID & \texttt{BG-THM-037} \\
Status & \texttt{proved} \\
Lean declaration & \path{ProjectKoios.BandGauge.Geometry.admissibleSetEDist_pos_iff_disjoint_of_compact} \\
Owning module & \path{ProjectKoios/BandGauge/Geometry/AdmissibleSets.lean} \\
Module SHA-256 & \texttt{e830b1b940ced88bc86b169261673337126583ee09e9ad9fb30ff580e9160fd8} \\
\bottomrule\end{tabularx}
\section{Exact mathematical statement}
For nonempty compact subsets $A,B$ of a metric space,
\[
  0<\delta_{\rm e}(A,B) \quad\Longleftrightarrow\quad A\cap B=\varnothing.
\]
Lean states the right side as \texttt{Disjoint A B}.
\section{Provenance and statement boundary}
The result closes the finite separation direction selected by the Formal proof companion
at merge \texttt{1078380fb608269a8094c7e86776a1c0d826333a}. It is classical compact
metric geometry specialized to the repository's extended distance; no novelty claim is
made.
\section{Proof argument}
If the distance is positive, a common point would force it to zero, so the sets are
disjoint. Conversely, nonempty compact disjoint sets attain their distance at distinct
points; metric separation of those points makes the attained distance strictly positive.
The proof reuses the existing lower-bound/disjointness theorem and compact positive
attainment theorem rather than introducing a new distance convention.
\section{Lean realization}
\begin{lstlisting}
theorem admissibleSetEDist_pos_iff_disjoint_of_compact
    {A B : Set X} (hA : IsCompact A) (hB : IsCompact B)
    (hneA : A.Nonempty) (hneB : B.Nonempty) :
    0 < admissibleSetEDist A B ↔ Disjoint A B
\end{lstlisting}
\section{Trusted base and review}
Lean \texttt{4.35.0}, Mathlib
\texttt{56a005855ea17e0b0958f73ec801ae140b0a1679}:
\begin{lstlisting}
depends on axioms: [propext, Classical.choice, Quot.sound]
\end{lstlisting}
No admissions or repository-owned axioms occur. AI-assisted and kernel-checked; no
independent expert review is recorded.
\section{Scientific and evidence boundaries}
The theorem needs exact compactness and nonemptiness. It does not certify a numerical
lower bound, authenticate data, establish external inventory coverage, or equate positive
tolerance-level map compatibility with intersection.
\end{document}
